\documentclass{homework}
\author{Tashfeen, Ahmad}
\class{CSCI 2114: Data Structures}
\title{Homework 6}
\address{Petree College of Arts \& Sciences, Computer Science, OCU}

\newcommand\providecode[1]{%
  \href{https://tashfeen.org/raw/share/ds/#1}{\texttt{#1}}
}

\begin{document} \maketitle

\question Head over to the \url{https://www.gutenberg.org} and find an
interesting book that is available in the text format. Perform word
frequency analysis on the text of the book using the Java's built-in
\href{https://docs.oracle.com/javase/8/docs/api/java/util/HashMap.html}{HashMap}.

Print out the top ten words with their respective frequencies in
descending order. The words should be at least six letters. For
example, a sample output for the George Orwell's controversial novel
\textit{1984} is as follows,

\begin{verbatim}
|   winston | 282|
|   thought | 158|
|    seemed | 133|
|    moment | 131|
|    always | 122|
|    obrien | 119|
|    almost | 106|
|   another | 105|
|    though | 104|
|    before | 103|
\end{verbatim}

Place the code for this question in a file called \texttt{Frequency.java}.

\question We used the following hash function in class to implement a
HashMap/Dictionary class in Java. Consider the ordering,
\[
  f = \{\text{(\textvisiblespace, 0), (a, 1), (b, 2), (c, 3),
    (d, 4),}\hdots, \text{(z, 26)}\}
\]
Then the hash function is as follows for a string $s$ of length $n$
consisting only of the lower-case letters and spaces,
\[
  h(s) = \sum_{i=0}^{n-1}(f(s_{i})\times 27^i) \mod m
\]
For example,
\begin{align*}
  h(\text{\textvisiblespace})  & =(0 \times 27^0)=0 \mod m                  \\
  h(\text{a})                  & =(1 \times 27^0)=1 \mod m                  \\
  h(\text{b})                  & =(2 \times 27^0)=2 \mod m                  \\
  h(\text{a\textvisiblespace}) & =(1 \times 27^1)+(0 \times 27^0)=27 \mod m
\end{align*}

In class we let $m=(2^{31}-1)/2$. Do you recall why?

Download the professor's implementation of the above hash function in
Java from \providecode{Base27.java}. Also download the list of
\href{https://raw.githubusercontent.com/dwyl/english-words/master/words_alpha.txt}
{370,105 unique English words}. Write a program that finds all the
hash collisions among these words for,
\begin{align*}
  m_1 & = (2^{31}-1)/2 = \texttt{Integer.MAX\_VALUE / 2} \\
  m_2 & = 2^{31}-1 = \texttt{Integer.MAX\_VALUE}         \\
  m_3 & = 2^{63}-1 = \texttt{Long.MAX\_VALUE}            \\
\end{align*}

For each hash collision corresponding to each $m_i$ give the hash
value and the collided words.

Hint: There are 89 hash collisions when
$m=(2^{31}-1)/2 = \texttt{Integer.MAX\_VALUE / 2}$ with one of them
being $h(\text{agonizes}) = 510264784 = h(\text{trioxides})$.  Use the
main method of \texttt{Base27.java} for this question.

\question Read Cormen, et al. \textit{Introduction to Algorithms}
($3^\text{rd}$ Edition) pages 647--652 and pages 658--659. You'll be
implementing the Dijkstra's and the Bellman-Ford\footnote{Also known
  as the ``distance vector routing'' algorithm in the networks field.}
algorithms. Download the files containing the adjacency matrices,

\begin{enumerate}
  \item \providecode{dijkstra1.txt}
  \item \providecode{dijkstra2.txt}
  \item \providecode{bellmanford1.txt}
  \item \providecode{bellmanford2.txt}
\end{enumerate}

The above matrices define edge weights between vertices that should be
named with labels in the list,
\[
  \text{a, b, c,} \hdots, \text{z}
\]
Taking as many labels as the rows/columns.

Write a Java class \texttt{Graph.java} that implements both the
Dijkstra's and the Bellman-Ford algorithms. Use your Dijkstra's
implementation to solve the files \texttt{dijkstra1.txt} and
\texttt{dijkstra2.txt}. Similarly, use your Bellman-Ford algorithm to
solve the files \texttt{bellmanford1.txt} and
\texttt{bellmanford2.txt}. For each graph and algorithm, report the
minimum distance with which each vertex can be reached and their
respective parent vertex. E. g, the output of Dijkstra's algorithm on
\texttt{dijkstra1.txt} should look like,

\begin{verbatim}
a: 0 via a
b: 8 via d
c: 9 via b
d: 5 via a
e: 7 via d
\end{verbatim}

You may verify the above by manually drawing the graph on paper and
running Dijkstra on it by hand.

\question Can you run the Bellman-Ford algorithm on
\texttt{dijkstra1.txt} and \texttt{dijkstra2.txt} to obtain the same
result as you did by running Dijkstra's algorithm? Similarly, can you
run Dijkstra's algorithm on \texttt{bellmanford1.txt} and
\texttt{bellmanford2.txt} to obtain the same results as you did by
running Bellman-Ford algorithm?

\question On a graph of exclusively positive edge weights, which
algorithm is a better choice?

\question Watch \href{https://www.youtube.com/watch?v=CgW0HPHqFE8}{A*
  (A-Star) Pathfinding Algorithm Visualization on a Real Map} on
YouTube to see how a path finding algorithm known as the A$^*$ looks
in execution on a real world graph. Note that A$^*$ is essentially the
Dijkstra's algorithm with addition of a ``heuristic'' function.

\end{document}
